\documentclass[11pt]{article}
\usepackage{amsmath,amssymb,amsthm,booktabs}
\usepackage[margin=1in]{geometry}
\newtheorem{observation}{Observation}
\title{The Proof Landscape of the Riemann Hypothesis:\\A Computational Map (Session 31)}
\author{8 iterations $\cdot$ 15 scripts $\cdot$ March 31, 2026}
\date{}
\begin{document}
\maketitle

\begin{abstract}
Starting from a proved local stability result (curvature $+1/\varepsilon^2$ at each zero), we systematically tested every computationally accessible approach to promoting local to global stability. Eight approaches were killed; four survive. The central discovery is that RH asserts a \emph{critical phenomenon}: the de Bruijn--Newman constant $\Lambda = 0$ is a phase transition, with zeros maximally close (GUE repulsion) without colliding. The obstruction to proof is the \emph{finite-to-infinite transition}: every approach works for finitely many primes/zeros but develops $O(1)$ errors in the infinite limit. The function field analogue (where RH is proved) avoids this because its zeta function is a polynomial.
\end{abstract}

\section{Approaches Killed}
\begin{center}
\begin{tabular}{rlp{7.5cm}}
\toprule
\# & Approach & Mechanism of death \\
\midrule
1 & Energy convexity & $d^2E/d\varepsilon^2 < 0$ at intermediate $\varepsilon$ \\
2 & Gamma confinement & Stirling terms cancel between conjugate pair \\
3 & Spacing lower bound & GUE: $\delta_{\min} \sim N^{-1/3} \to 0$ \\
4 & Information-theoretic & Explicit formula is a tautology \\
5 & Li criterion computation & Detection requires $n \sim \gamma^2/\delta$ \\
6 & Anti-alignment & Max interference at $\sigma \approx 0$, not $1/2$ \\
7 & Moment constraint (100\%) & Double/sum ratio $= 1.003$ (too soft) \\
8 & Spectral Jacobi matrix & Trivially self-adjoint by construction \\
\bottomrule
\end{tabular}
\end{center}

\section{Approaches Surviving}
\begin{center}
\begin{tabular}{rp{4cm}p{5cm}l}
\toprule
\# & Approach & Status & Wall \\
\midrule
1 & Connes $Q_W$ positivity & Confirmed as Castelnuovo analogue & $O(1)$ leakage \\
2 & Levinson--Conrey mollifiers & 40\% proved & Mean-value theorems \\
3 & GUE universality & Data confirms match & Open problem \\
4 & Rodgers--Tao dual & Unexplored & Technical \\
\bottomrule
\end{tabular}
\end{center}

\section{The Three Key Discoveries}

\begin{observation}[RH is a critical phenomenon]
$\Lambda = 0$ means the zero configuration is at the exact phase transition. GUE predicts $\delta_{\min} \sim N^{-1/3}$, giving collision time $t_c \sim N^{-2/3} \to 0^+$. The infimum of all collision times is $\Lambda = 0$ exactly. There is no safety margin.
\end{observation}

\begin{observation}[The finite-to-infinite obstruction]
Function field RH (proved): zeta is a polynomial, proof uses algebraic geometry of a finite-dimensional cohomology. Number field RH (open): zeta is transcendental, proof would require analytic control of an infinite-dimensional operator. Every approach (Connes, Levinson--Conrey, Li, moments) works for finite truncations but develops $O(1)$ errors in the infinite limit.
\end{observation}

\begin{observation}[Euler product mechanism]
The partial Euler product $\prod_{p \leq 997}(1-p^{-s})^{-1}$ equals $0.061$ at a zero (far from zero---the infinite tail does the work) but $1.006 \times \zeta(s)$ off the critical line (essentially converged). The prime conspiracy (predominantly negative $\cos(t\log p)$ at zeros) creates zeros, but through an infinite coordination that no finite truncation captures.
\end{observation}

\section{The Proof Map}

All roads lead to the same wall:

\medskip
\begin{center}
\textbf{Euler product} $\longrightarrow$ \textbf{GUE universality} $\longrightarrow$ \textbf{$\Lambda = 0$} $\longrightarrow$ \textbf{RH}
\end{center}
\medskip

\noindent Each arrow is a major open problem. The Connes $Q_W$ positivity is the most direct attack: it is the number field analogue of the Castelnuovo--Weil inequality that proves the function field case.

\section{Quantitative Summary}
15 scripts, 500 zeros, $\sim$50 computational experiments:
\begin{itemize}
\item GUE fit: $\chi^2 = 0.63$ (vs Poisson $10.6$, ratio $16.9\times$)
\item Sign changes: $150/150$ confirmed, $S(T) = N(T)$
\item Min spacing ratio: $\delta_{\min}/\delta_{\mathrm{avg}} \geq 0.191$ (500 zeros)
\item Euler partial product at zero: $|{\textstyle\prod_{p\leq 997}}| = 0.061$
\item Moment double/sum: $1.003$ (indistinguishable)
\item Codimension $C_R$--$C_I$ gap: $0.6$--$1.5$ (shrinking, slope $-0.018$)
\end{itemize}

\end{document}
